\section{Metodoloxía e planificación}

\begin{frame}{Desenvolvemento iterativo e incremental}
  \begin{columns}[c]
    \begin{column}{0.42\textwidth}
      \small
      \begin{itemize}
        \item Once iteracións, onde a maioría corresponde cun
        compoñente do hardware da Game Boy.

        \item A orde vén imposta polo acoplamento do hardware.

        \item O desenvolvemento dos compoñenetes segue unha estrutura de estudo da arquitectura, implementación
          e validación.

        \item Git e GitHub, empregando \textit{conventional commits} e \textit{conventional branch}
        para manter un historial de cambios limpo.

      \end{itemize}
    \end{column}
    \begin{column}{0.56\textwidth}
      \begin{center}
      \begin{ganttchart}[
          hgrid, vgrid,
          x unit=0.38cm,
          y unit title=0.40cm,
          y unit chart=0.32cm,
          title height=1,
          title label font=\tiny\bfseries,
          bar/.style={fill=udcpink!75, draw=udcpink!90},
          bar height=0.62,
          bar label font=\tiny,
        ]{1}{18}
        \gantttitle{2025}{6}\gantttitle{2026}{12} \\
        \gantttitle{O}{2}\gantttitle{N}{2}\gantttitle{D}{2}
        \gantttitle{X}{2}\gantttitle{F}{2}\gantttitle{M}{2}
        \gantttitle{A}{2}\gantttitle{M}{2}\gantttitle{X}{2} \\
        \ganttbar{It. 1}{1}{1} \\
        \ganttbar{It. 2}{1}{2} \\
        \ganttbar{It. 3}{2}{5} \\
        \ganttbar{It. 4}{4}{6}\ganttbar{It. 4}{9}{9} \\
        \ganttbar{It. 5}{3}{6} \\
        \ganttbar{It. 6}{9}{10} \\
        \ganttbar{It. 7}{10}{11} \\
        \ganttbar{It. 8}{11}{13} \\
        \ganttbar{It. 9}{12}{13} \\
        \ganttbar{It. 10}{14}{15} \\
        \ganttbar{It. 11}{15}{16} \\
        \ganttbar[bar/.append style={fill=udcgray!50, draw=udcgray}]{Doc.}{15}{18}
      \end{ganttchart}\\[0.4em]
      \fonte{Diagrama de Gantt do proxecto.}
      \end{center}
    \end{column}
  \end{columns}
\end{frame}
